% SOURCE_AUTHORITY: sga5_fr_workpass.tex, lines 5523--5542 (LF-normalized unit).
% SOURCE_SHA256: 791F4EFFC5E02832D5D77ED03518C8156D6F07E4C8238B03545DB93D883FBB28
\section*{6. Correspondencias equivariantes y divisibilidad de los términos locales}

\subsection*{6.1. Notaciones}

En este número, $S$ designa el espectro de un cuerpo de característica $p$, y $\Lambda$ un anillo conmutativo noetheriano anulado por un entero no divisible por $p$. Se supondrá que los $S$-esquemas considerados son separados y de tipo finito. Las $\Lambda$-álgebras consideradas serán asociativas, unitarias y de tipo finito, pero no necesariamente conmutativas. En la mayoría de los casos, y en particular para todo lo relativo a los ``términos locales de Lefschetz--Verdier no conmutativos'', se supondrá que son planas; en las aplicaciones serán álgebras de grupos finitos.

Si $X$ es un $S$-esquema y $A$ una $\Lambda$-álgebra, se denotará por $D({}_AX)$ (resp. $D(X_A)$) la categoría derivada de la categoría de haces de $A$-módulos izquierdos (resp. derechos) sobre $X$; para $A$ conmutativa, se escribirá también $D(X,A)$ en lugar de $D({}_AX)$, y $D(X)=D(X,\Lambda)$. Se denotará por $D_c(-)$ (resp. $D_{\mathrm{ctf}}(-)$) la subcategoría plena formada por los complejos con cohomología constructible (resp. de dimensión Tor finita y con cohomología constructible). Si $A$, $B$ son $\Lambda$-álgebras, se denotará por $D({}_AX_B)$ la categoría derivada de la categoría de haces de $(A,B)$-bimódulos, izquierdos sobre $A$ y derechos sobre $B$. Por último, se utilizarán las notaciones $A^o$, $A^e$, $A_{\natural}$ de 5.1.

\subsection*{6.2}
El formalismo de (III 1, 2, 3) se extiende sin dificultad al caso ``no conmutativo''. Indiquemos brevemente los puntos principales de esta generalización.

Sean $X$ un $S$-esquema, y $A$, $B$ unas $\Lambda$-álgebras. Si $L\in\Ob D_{\mathrm{ctf}}(X_A)$, $M\in\Ob D({}_AX_B)$, siendo $M$ de dimensión Tor finita y con cohomología constructible como objeto de $D(X_B)$, entonces
\[
L\lten_A M\in\Ob D(X_B)_{\mathrm{ctf}}
\]
(la comprobación es inmediata). Por otra parte, de (SGA $4^{1/2}$ Finitude 1.6) se deduce que, para $M$ como antes y $L\in\Ob D_{\mathrm{ctf}}({}_AX)$, se tiene
\[
\uRHom_A(L,M)\in\Ob D_{\mathrm{ctf}}(X_B).
\]
De (loc. cit.) resulta asimismo que, si $f$ es un $S$-morfismo, $D_{\mathrm{ctf}}({}_A-)$ es estable bajo las operaciones $(f^*,f_*,f_!,f^!)$ (siendo trivial el caso de $f^*$).
